% Pregunta 7. Desequilibrios, riesgos y recomendaciones (bloque VI) · peso alto
\clearpage
\section*{Pregunta 7. Desequilibrios, riesgos y recomendaciones}

\begin{enunciado}
\textbf{7.} Identifique los principales \textbf{desequilibrios y vulnerabilidades} de la economía española y los riesgos al alza y a la baja que rodean las previsiones para 2026 y 2027. Teniendo en cuenta las tendencias demográficas, ¿qué \textbf{reformas estructurales y políticas de oferta} recomendaría para sostener el crecimiento a medio y largo plazo?

\smallskip
\textbf{Peso alto} · \tiempo{15} · guía: VI.1 (fortalezas, desequilibrios y medidas; naturaleza, origen y gravedad), VI.2 (recomendaciones por desequilibrio), V.5 (políticas de oferta), D.8 (matriz), B.1 (demografía) · todo el cuadernillo.
\end{enunciado}

\begin{examen}
\textit{España afronta el final de la década con fortalezas que no tenía en 2007 (superávit exterior, sector privado desapalancado, banca rentable, empleo en máximos) y con 3 desequilibrios de fondo, por este orden: una productividad estancada, unas cuentas públicas con déficit estructural y deuda alta, y un paro estructural elevado; a ellos se suma, en 2026, un choque de precios que puede volverse persistente. El envejecimiento hace más urgente corregirlos, porque el crecimiento no podrá apoyarse indefinidamente en sumar trabajadores.}

\rotulo{a) Desequilibrios y vulnerabilidades}
Las \textbf{fortalezas} son claras: un crecimiento superior al de la zona del euro, la tasa de empleo más alta de la serie, un superávit exterior que se mantiene desde 2012, la deuda privada en su mínimo (hogares, 42,9~\%; empresas, 62,0~\% del PIB), una banca rentable y una prima de riesgo de 45~pb. Los desequilibrios, ordenados por su importancia:

\textbf{1.}~\textbf{Productividad estancada}: el PIB por hora crece un 0,33~\% anual desde 2019, el gasto en I+D es dos tercios del europeo y la inversión en maquinaria apenas crece. Es estructural y persistente, y es el más importante porque \uline{solo la productividad eleva la renta por habitante}, que desde 2007 crece menos que en la zona del euro.


\textbf{2.}~\textbf{Finanzas públicas}: un déficit estructural del 2,3~\% con la economía en su potencial, que las previsiones no reducen, y una deuda del 100,5~\%, que ha bajado gracias al crecimiento nominal. Es un problema de \emph{stock}, persistente, que se agravará con el envejecimiento y con la normalización del coste de la deuda.


\textbf{3.}~\textbf{Paro estructural}: el 10,5~\%, el doble que en la zona del euro, con un paro juvenil del 24,9~\% y muy pocas vacantes: un problema de emparejamiento y de cualificación más que de demanda.


\textbf{4.}~\textbf{Inflación y costes}: coyuntural, pero con riesgo de enquistarse; la inflación del 5,0~\% y unos costes laborales que crecen más que en la zona erosionan la competitividad ganada.

Los 2 primeros son problemas de \emph{stock} y de oferta, que no se corrigen solos con el ciclo; el tercero es en parte la histéresis de las crisis pasadas; el cuarto es de flujo, pero los problemas de flujo que persisten acaban siendo de \emph{stock}, y por eso conviene atajarlo pronto. Y 3 \textbf{vulnerabilidades}: la dependencia energética (el déficit de bienes llegó al 4,4~\% del PIB en 2022 y el superávit corriente se reducirá en 2026), la vivienda (precios a 2 dígitos con una oferta que no responde) y una posición exterior todavía por debajo del umbral europeo, aunque cada vez menos preocupante.

\rotulo{b) Riesgos de las previsiones}
Las previsiones del FMI (PIB del 2,5 y el 1,8~\%; inflación del 3,7 y el 2,7~\%) ya incorporan el choque del petróleo. Los riesgos para el \textbf{crecimiento} son, en 2026, equilibrados: el ritmo del primer semestre (0,6-0,7~\% trimestral) es coherente con la previsión, pero el petróleo, por encima de los 110 dólares en septiembre, resta renta real a los hogares y márgenes a las empresas, y empeora la balanza de bienes (en el choque de 2022 el déficit de bienes pasó del 1,7 al 4,4~\% del PIB). En 2027 dominan los riesgos \textbf{a la baja}: el encarecimiento del crédito (el Euríbor ha subido casi 1 punto y los tipos de los préstamos se reprecian deprisa), la debilidad de los socios por los aranceles, la consolidación fiscal pendiente y un posible ajuste de la vivienda. \textbf{Al alza}: la llegada de población, el turismo, un ahorro de los hogares todavía alto (11,9~\%) que da margen al consumo, y la inversión en intangibles. Para la \textbf{inflación}, los riesgos son al alza: la previsión anual (3,7~\%) exige que el último trimestre baje al 4~\% desde el 5,0~\% de septiembre, y la subyacente no cede.

\rotulo{c) Reformas estructurales y políticas de oferta}
\textbf{A corto plazo}, con una inflación del 5,0~\% y una subyacente del 3,5~\% (gráfico~6), la prioridad sería que el choque energético no se convierta en inflación persistente: una política fiscal que no añada demanda (ayudas, si las hay, temporales y focalizadas en los hogares vulnerables y que mantengan el incentivo a ahorrar energía), un pacto de rentas y una política macroprudencial preparada para frenar el crédito hipotecario si se acelera.\footnote{El FMI recomienda retirar las medidas energéticas al expirar y, en un escenario grave, sustituirlas por ayudas focalizadas que no alteren los precios \fmi{17}.} La política monetaria es común y no puede ajustarse a España.

\textbf{A medio y largo plazo}, el punto de partida es la \textbf{demografía}: la tasa de dependencia de los mayores sube del 24,0 al 31,2~\% entre 2007 y 2025 (cuadro~3), y el peso de la población en edad de trabajar ya no suma al crecimiento (cuadro~2~bis), que se ha apoyado en la inmigración. Para sostenerlo, haría falta que creciera la productividad y que trabajara más gente. Las recomendaciones, ordenadas por la importancia del desequilibrio que corrigen, con su objetivo y su riesgo:
\begin{itemize}
\item \textbf{Productividad}: con un PIB por hora que crece un 0,33~\% anual desde 2019 (cuadro~2~bis) y un gasto en I+D del 1,50~\% del PIB, frente al 2,26~\% de la UE (cuadro~6), convendría reforzar el capital humano (menos abandono escolar, formación profesional dual, recualificación), la I+D (incentivos fiscales, colaboración entre universidades y empresas, capital riesgo) y un entorno que permita a las empresas crecer (menos trabas regulatorias ligadas al tamaño, unidad de mercado, competencia en los servicios). Objetivo: un crecimiento más intensivo. Riesgo: sus efectos tardan años.\footnote{El FMI recomienda una reforma de la universidad que reduzca el desajuste de cualificaciones, suprimir umbrales fiscales y regulatorios ligados al tamaño, avanzar en la unidad de mercado («Régimen 20») y en la integración de los mercados de capitales europeos, y rediseñar el crédito fiscal a la I+D para que llegue a las empresas jóvenes \fmi{30-31}.}
\item \textbf{Mercado de trabajo}: con un paro del 10,5~\%, un paro juvenil del 24,9~\% y una tasa de vacantes de solo el 0,9~\% (cuadro~3, gráfico~3), serían útiles unas políticas activas de empleo focalizadas en jóvenes y parados de larga duración y evaluadas, una mejor intermediación de los servicios públicos de empleo, consolidar la menor temporalidad (15,4~\%) y salarios ligados a la productividad. Objetivo: bajar el paro estructural. Riesgo: la eficacia de las políticas activas depende de su diseño y su evaluación.\footnote{El FMI recomienda financiar los servicios públicos de empleo según su tasa de colocación, cotizaciones más altas para las empresas que despiden y contratan con más frecuencia, ampliar la compatibilidad del subsidio con el empleo y no subir el salario mínimo por encima del 60~\% del salario medio neto \fmi{28-29}.}
\item \textbf{Demografía y pensiones}: con una población extranjera que pasa del 10,3 al 14,1~\% del total desde 2019 y una tasa de dependencia al alza (cuadro~3), convendría integrar bien la inmigración (reconocimiento de cualificaciones, formación), elevar la actividad de los mayores y de las mujeres (jubilación flexible, servicios de cuidados) y asegurar la sostenibilidad de las pensiones ligándolas a la evolución demográfica. Objetivo: que la población activa siga creciendo aunque aumenten las jubilaciones.\footnote{El FMI propone, entre otras opciones, alargar el número de años cotizados que se usan para calcular la pensión \fmi{19}.}
\item \textbf{Consolidación fiscal gradual}: con un déficit estructural del 2,3~\% que las previsiones no reducen y una deuda del 100,5~\% (cuadro~9), unas décimas del PIB al año hasta llevar el déficit estructural por debajo del 1,5~\% y lograr superávits primarios, con revisiones del gasto, menos beneficios fiscales y bases más amplias, protegiendo la inversión pública. Riesgo: frenar la demanda si se hace deprisa.\footnote{El FMI recomienda 0,5 puntos al año entre 2026 y 2030, con medidas adicionales por un 1,5~\% del PIB, sobre todo armonizando los tipos reducidos del IVA \fmi{17-18}.}
\item \textbf{Vivienda}: los precios suben un 12,7~\% en 2025 (un 9,7~\% en términos reales), por encima del umbral del 9~\% del procedimiento de desequilibrios, pero, a diferencia de 2007, sin que los impulse el crédito: la deuda de los hogares está en mínimos de la serie (42,9~\% del PIB; cuadro~4, gráfico~4). Al mismo tiempo, la población y el empleo crecen con fuerza (un 0,7 y un 2,6~\% en 2025; cuadro~3) y la inversión en vivienda apenas aumenta (un 6,0~\% desde 2019, frente a un 13,4~\% del total de la inversión; cuadro~1). Todo apunta a un desajuste entre una demanda que crece y una oferta que, al menos a corto plazo, no la absorbe. Como el precio depende de ambos lados del mercado, convendría actuar sobre los 2: por el lado de la oferta, facilitando que se construya más y ampliando el parque de alquiler asequible; por el lado de la demanda, moderando la demanda no residencial o especulativa (por ejemplo, con la fiscalidad) y vigilando el crédito hipotecario con instrumentos macroprudenciales por si se acelera. Serían menos aconsejables las ayudas generales a la compra, que con una oferta poco flexible tenderían a trasladarse a los precios. La vivienda cara podría estar dificultando también la movilidad geográfica y el emparejamiento en el mercado de trabajo (las vacantes están en su máximo desde 2011 con un paro todavía del 10,5~\%; cuadro~3, gráfico~3).\footnote{El FMI recomienda acelerar el desarrollo urbanístico, simplificar las licencias, retomar la reforma de la ley del suelo y que el control de alquileres sea temporal y esté condicionado a aumentar la oferta \fmi{29, 81}.}
\item \textbf{Energía}: con un déficit de bienes que llegó al 4,4~\% del PIB en el choque de 2022 y que las previsiones sitúan de nuevo en el 4,3-4,5~\% en 2026-2027 (cuadro~7), convendría diversificar las fuentes, acelerar las renovables y reforzar las interconexiones y la eficiencia, para reducir la exposición a choques como los de 2022 y 2026.\footnote{El FMI señala que el peso creciente de las renovables amortigua el traslado del precio del gas a la electricidad \fmi{11}.}
\end{itemize}
En conclusión, España tiene hoy un margen que no tuvo en 2008 (superávit exterior y deuda privada baja; cuadros~4 y~7) para \uline{aprovechar la expansión y corregir los desequilibrios públicos y de productividad}; cómo lo haga decidirá si su crecimiento se traduce en renta por habitante.
\end{examen}

\begin{complemento}{7.1 · La matriz de fortalezas y debilidades, rellena}
\que{Qué se aplica.} Guía, D.8: la matriz se va rellenando durante el examen en una hoja aparte y sirve para 3 cosas: anticipar el diagnóstico en la pregunta~1, no contradecirse y contestar la última pregunta. VI.1: \emph{qué va bien, qué hay que cambiar (2 o 3 desequilibrios ordenados) y cómo}; las 3 preguntas del FMI (naturaleza, origen y gravedad) \parencite{FMI2013CaseStudy}.

\begin{tablacaja}
\centering\scriptsize
\renewcommand{\arraystretch}{0.98}\setlength{\tabcolsep}{4pt}
\caption{Matriz de fortalezas y debilidades de España (examen de ejemplo)}\label{tab:matriz-espana}
\begin{tabular}{>{\raggedright\arraybackslash}p{2.0cm}>{\raggedright\arraybackslash}p{3.3cm}>{\raggedright\arraybackslash}p{3.6cm}>{\centering\arraybackslash}p{0.7cm}>{\centering\arraybackslash}p{0.8cm}>{\raggedright\arraybackslash}p{3.4cm}}
\toprule
\rowcolor{tceeazulpalido}
\textbf{Bloque} & \textbf{\textcolor{tceeverde}{Fortaleza}} & \textbf{\textcolor{tceerojo}{Debilidad}} & \textbf{Peso} & \textbf{Tend.} & \textbf{Recomendación} \\
\midrule
Perfil cíclico & PIB +10,0~\% en 2019-2025 (zona: +6,8~\%) & Crecimiento extensivo: PIB por habitante +5,0~\% & ** & $\nearrow$ & Políticas de oferta (productividad) \\
Trabajo & Empleo 67,0~\%; temporalidad 15,4~\% & Paro 10,5~\%; juvenil 24,9~\%; PIB por hora +0,33~\%/año & *** & $\nearrow$ & Políticas activas, emparejamiento, capital humano \\
Capital y crédito & Deuda privada en mínimos; intangibles +47~\% & Maquinaria +2,7~\%; vivienda +12,7~\% & * & $\rightarrow$ & Vivienda: oferta y demanda; incentivos a invertir \\
Precios & Salida rápida de 2022-2023 & IPCA 5,0~\% (zona 3,8); CLU +4,4~\% & ** & $\searrow$ & Pacto de rentas; fiscal que no empuje \\
No precio & Cuota mundial estable; servicios & I+D 1,50~\% (UE 2,26) & * & $\rightarrow$ & I+D, tamaño, diversificar mercados \\
Exterior & Superávit desde 2012; NENDI $-9{,}7~\%$ & PIIN $-42{,}7~\%$; energía & * & $\nearrow$ & Política energética \\
Monetaria & Transmisión rápida; diferencial 45~pb & Tipo real más bajo en España & * & $\rightarrow$ & (común; compensa la fiscal) \\
Banca & ROE 14,2~\%; dudosos 2,4~\% & CET1 13,9~\% (UE 16,4) & * & $\nearrow$ & Colchones; límites LTV/LTI si hace falta \\
Fiscal & Saldo sin intereses equilibrado; deuda a la baja & Déficit estructural 2,3~\%; deuda 100,5~\% & *** & $\nearrow$ & Consolidación gradual con la senda de gasto \\
Demografía & Inmigración (extranjeros 14,1~\%) & Dependencia 31,2~\% & ** & $\searrow$ & Integración, actividad, pensiones \\
\bottomrule
\end{tabular}
\fuente{Elaboración propia con el cuadernillo; plantilla de la guía (D.8). Peso: importancia del desequilibrio (*** máxima). Tendencia: $\nearrow$ mejora, $\rightarrow$ estable, $\searrow$ empeora.}
\end{tablacaja}
\end{complemento}

\begin{complemento}{7.2 · Los umbrales oficiales: qué se enciende en España}

\begin{tablacaja}
\centering\scriptsize
\caption{España frente a los umbrales del procedimiento de desequilibrios y del Tratado}\label{tab:pdm}
\begin{tabular}{>{\raggedright\arraybackslash}p{4.4cm}>{\centering\arraybackslash}p{2.6cm}>{\centering\arraybackslash}p{3.2cm}>{\centering\arraybackslash}p{2.4cm}}
\toprule
\rowcolor{tceeazulpalido}
\textbf{Indicador} & \textbf{Umbral} & \textbf{España} & \textbf{¿Alerta?} \\
\midrule
Saldo corriente, media de 3 años & $-4$ / $+6~\%$ & $(2{,}7+3{,}3+3{,}2)/3=3{,}1$ & \textcolor{tceeverde}{no} \\
PIIN & $-35~\%$ & $-42{,}7~\%$ & \textcolor{tceerojo}{sí} (mejorando) \\
Tipo de cambio efectivo real, 3 años & $\pm3~\%$ & $-0{,}9~\%$ (2022-2025) & \textcolor{tceeverde}{no} \\
CLU nominales, 3 años & $+9~\%$ & $\approx+13{,}9~\%$ (2023-2025)* & \textcolor{tceerojo}{sí} \\
Precio de la vivienda (nominal) & $+9~\%$ & $+12{,}7~\%$ & \textcolor{tceerojo}{sí} \\
Deuda de hogares / de empresas & 55 / 85~\% & 42,9 / 62,0~\% & \textcolor{tceeverde}{no} \\
Tasa de paro & 10~\% & 10,5~\% & \textcolor{tceerojo}{sí} (por poco) \\
Saldo público (Tratado) & $-3~\%$ & $-2{,}4~\%$ & \textcolor{tceeverde}{no} \\
Deuda pública (Tratado) & 60~\% & 100,5~\% & \textcolor{tceerojo}{sí} \\
\bottomrule
\end{tabular}
\fuente{Umbrales: \textcite{CE2025AMR2026,TFUE2012Protocolo12} (guía, A.1); datos: cuadros 3, 4, 6, 7 y 9. (*) El umbral se refiere a los CLU por hora; con los CLU por persona del cuadernillo, $1{,}059\times1{,}030\times1{,}044-1\approx13{,}9~\%$, una aproximación. Qué hay que ver: las alertas están en los costes, la vivienda, el paro y la deuda pública: exactamente los desequilibrios de la respuesta. Son señales, no obligaciones (guía, A).}
\end{tablacaja}
\begin{figuracaja}
\caption{Los umbrales oficiales de un vistazo, con la posición de España}\label{fig:umbrales}
\grafico[0.66\linewidth]{s09-umbrales-espana}
\fuente{Infografía de la guía (A, umbrales del procedimiento de desequilibrios y del Tratado) con el último dato de España del cuadernillo. Qué hay que ver: los rombos en rojo (PIIN, paro, CLU, vivienda, deuda pública) son los desequilibrios de la respuesta; los demás, en verde.}
\end{figuracaja}
\end{complemento}

\begin{complemento}{7.3 · Riesgos, recomendaciones que responden a los datos y qué recortar}
\que{Del dato a la recomendación} (guía, VI.1: «puntúan las recomendaciones que responden a los datos del cuadro»): cada medida de la respuesta sale de un indicio (gráfico~\ref{fig:dato-recomendacion}). Ejemplo: \emph{PIB por hora +0,33~\% e I+D del 1,50~\%} $\to$ productividad estancada $\to$ I+D y tamaño empresarial $\to$ objetivo: crecimiento intensivo $\to$ riesgo: tarda.

\begin{figuracaja}
\caption{Del dato a la recomendación (infografía de la guía)}\label{fig:dato-recomendacion}
\grafico[0.56\linewidth]{esquema-dato-recomendacion}
\fuente{Guía del dictamen, VI.1. Qué hay que ver: siempre el mismo camino, del indicio al desequilibrio, a la medida, a su objetivo y a su riesgo.}
\end{figuracaja}

\begin{tablacaja}
\centering\scriptsize
\caption{Riesgos de las previsiones de 2026-2027 y su indicio en el cuadernillo}\label{tab:riesgos}
\begin{tabular}{>{\raggedright\arraybackslash}p{1.6cm}>{\raggedright\arraybackslash}p{5.2cm}>{\raggedright\arraybackslash}p{7.0cm}}
\toprule
\rowcolor{tceeazulpalido}
\textbf{Sentido} & \textbf{Riesgo} & \textbf{Indicio} \\
\midrule
A la baja & Petróleo caro y persistente & Brent de 62,5 a 114,2 dólares (G7); bienes $-4{,}4~\%$ del PIB en el choque de 2022 (C7) \\
 & Endurecimiento financiero & Depósito 2,50~\%; Euríbor +98~pb desde diciembre de 2025; hipotecas 3,04~\% (G12) \\
 & Aranceles y socios & Exportaciones a EE.~UU. $-11~\%$ (G9) \\
 & Menos transferencias de capital & Cuenta de capital: 1,3~\% del PIB en 2026 y 0,2~\% en 2027 (C7) \\
 & Vivienda & Precio +12,7~\% con poca inversión residencial (C1, C4) \\
Al alza & Población y turismo & Población +1,0~\% al año; viajes 4,3~\% del PIB (C3, C7) \\
 & Consumo & Ahorro de los hogares 11,9~\%, el doble que en 2007 (C4) \\
 & Ritmo de 2026 & 0,6-0,7~\% trimestral en el primer semestre (G1): $\approx$2,5~\% anual \\
Inflación (al alza) & Segunda ronda & Media de enero-septiembre 3,6~\% con septiembre en el 5,0; previsión 3,7~\% (C5); subyacente 3,5~\% (G6) \\
\bottomrule
\end{tabular}
\fuente{Cuadernillo (C: cuadro; G: gráfico).}
\end{tablacaja}

\que{Errores típicos} (guía, VI.1 y D.6): recomendaciones genéricas que valen para cualquier país («hacer reformas estructurales»); recetas tajantes o demasiado precisas; olvidar la demografía, que el enunciado pide expresamente; no ordenar los desequilibrios. \que{Alternativa de redacción} del cierre: «España tiene que aprovechar una expansión sin desequilibrios privados para corregir los públicos y elevar la productividad, antes de que el envejecimiento reduzca el crecimiento de la población activa». \que{Qué añadir con más tiempo.} El caso de España del FMI (guía, C.2): ajuste fiscal de unos 0,5 puntos al año entre 2026 y 2030, políticas activas y de vivienda, y reformas de productividad \parencite{FMI2026ArticuloIV}. \que{Qué recortar en el examen.} Las vulnerabilidades (1 frase) y 2 de las 6 recomendaciones (energía y vivienda, ya tratadas en las preguntas 1, 3, 4 y 5).
\end{complemento}

\begin{complemento}{7.4 · Contexto y fuentes: lo que el cuadernillo no trae}
\que{Riesgos que el cuadernillo no muestra.} El final del Mecanismo de Recuperación y Resiliencia en 2026 (Reglamento (UE) 2021/241), que ha financiado buena parte de la inversión en intangibles e infraestructuras, es un riesgo a la baja para 2027; el gasto en defensa (cláusula de escape de 2025, \cite{CE2025Defensa}), una presión sobre las cuentas públicas.

\que{El diagnóstico del FMI.} La consulta del artículo~IV de 2026 resume que España «siguió creciendo más que la zona del euro pese al impacto adverso de la guerra en Oriente Medio» y recomienda un ajuste fiscal gradual, políticas activas y de vivienda que mejoren el emparejamiento y reformas de productividad (capital humano, tamaño empresarial, I+D) \parencite{FMI2026ArticuloIV} (guía, C.2): las mismas líneas que salen de los datos en el texto de examen.

\que{Envejecimiento.} La Comisión proyecta un aumento del gasto ligado a la edad (pensiones, sanidad, cuidados) en la mayoría de los países de la UE \parencite{CE2024Envejecimiento}; con la tasa de dependencia de los mayores en el 31,2~\% y subiendo (cuadro~3), es el argumento de fondo para empezar ya la consolidación y las reformas de actividad.

\que{Las reformas ya hechas.} La reforma laboral de 2021 (Real Decreto-ley 32/2021, pregunta~2) es el ejemplo de una reforma cuyo efecto se ve en el cuadernillo (temporalidad del 25,4 al 15,4~\%): sirve para defender que las reformas de oferta funcionan, aunque sus efectos sobre la productividad tarden más en verse.
\end{complemento}

\begin{complemento}{7.5 · El diagnóstico del FMI (artículo~IV de 2026): riesgos y reformas}
\que{La matriz de riesgos} \fmi{52-56}. Con probabilidad alta: las tensiones geopolíticas (el conflicto de Oriente Medio es el principal riesgo a la baja), el proteccionismo, la volatilidad de las materias primas (España importa en torno al 70~\% de la energía que usa), una corrección brusca ligada a la inteligencia artificial, la subida de los tipos por la deuda pública mundial y los ciberataques. Riesgos internos: la fragmentación política (impacto alto), la ejecución de los fondos europeos (unos desembolsos del 1,2~\% del PIB en 2026) y el auge de la vivienda. Al alza: más turismo, más desahorro de los hogares, más inmigración (el potencial de 2030 sería 0,1 puntos mayor si se mantuviera la de 2025) y la inteligencia artificial (0,1-0,2 puntos más de productividad al año) \fmi{14}. \textbf{Coincide} con la tabla de riesgos de la solución y añade 2 que el cuadernillo no permite ver: la fragmentación política y la ejecución de los fondos.

\que{Las reformas} \fmi{28-31}. Mercado de trabajo (cuadro~2.5); vivienda: acelerar el desarrollo urbanístico, simplificar las licencias, recuperar la reforma de la ley del suelo y que el control de alquileres sea temporal y condicionado a aumentar la oferta (en Cataluña rebajó los alquileres un 4-6~\%, pero redujo los anuncios de alquiler entre un 15 y un 30~\%) \fmi{29, 81}; productividad: unidad de mercado («Régimen 20»), integración de los mercados de capitales europeos, reforma de la universidad para reducir el desajuste de cualificaciones, menos umbrales regulatorios ligados al tamaño y un crédito fiscal a la I+D más sencillo y dirigido a las empresas jóvenes, que podría elevar el crecimiento de la productividad total de los factores en un cuarto de punto \fmi{30-31}.

\textbf{Coincide} con las 6 líneas de la solución; \textbf{añade} detalle (el IVA y las pensiones en lo fiscal, el crédito fiscal a la I+D en la innovación) y una advertencia de economía política: con la fragmentación parlamentaria, el riesgo de no cumplir los planes es alto \fmi{34}.
\end{complemento}

\begin{complemento}{7.6 · La vivienda: ¿problema de oferta o de demanda? Las tijeras de \autor{Marshall}}
\que{Los 2 filos.} \autor{Marshall} comparó la oferta y la demanda con los 2 filos de unas tijeras: preguntarse cuál de los 2 determina el precio es como preguntarse qué filo corta el papel \parencite[libro~V, cap.~III]{Marshall1890}. Decir que la vivienda «es solo un problema de oferta» elige un filo sin pruebas; el cuadernillo muestra los 2: una demanda que crece con la población y el empleo y una inversión residencial que apenas responde (texto de examen, apartado~c).

\que{Qué filo conviene mover depende de las elasticidades.} A corto plazo el parque de viviendas es casi fijo (las viviendas terminadas están en un 14~\% del nivel de 2007 \fmi{47}). Con una oferta tan rígida:
\begin{itemize}
\item las medidas que \textbf{reducen la demanda} bajan la presión sobre el precio;
\item las que la \textbf{apoyan} sin ampliar la oferta se trasladan en buena parte al precio (se \emph{capitalizan}).
\end{itemize}
A largo plazo la oferta es más elástica, y ahí es donde rinden las medidas de oferta (suelo, licencias, seguridad jurídica), que tardan años.

\que{La evidencia sobre medidas de demanda}, con sus fuentes:
\begin{itemize}
\item \emph{Ayudas generales a la compra}: el programa británico \emph{Help to Buy} elevó los precios en Londres, donde la oferta está muy restringida, más que el valor de la subvención, sin efecto apreciable sobre la construcción \parencite{Carozzi2024HelpToBuy}.
\item \emph{Ayudas al alquiler}: en Francia, la ampliación de las ayudas a la vivienda de principios de los años 90 elevó los alquileres sin mejorar la calidad de las viviendas \parencite{Fack2006Ayudas}; según la reseña del Banco de Francia, la subida de los alquileres equivalió al 78~\% del importe de las ayudas \parencite{BdF2023AyudasVivienda}.
\item \emph{Fiscalidad de la vivienda vacía}: el impuesto francés sobre la vivienda vacía redujo su número en torno a un 13~\% entre 1997 y 2001, sobre todo la vacía de larga duración, y la mayoría pasó a ser vivienda principal \parencite{Segu2020Vacantes}.
\item \emph{Pisos turísticos}: en Barcelona, la actividad de Airbnb elevó los alquileres un 1,9~\% y los precios de compraventa un 4,6~\% en el barrio medio, y hasta un 7 y un 17~\% en los barrios con más actividad \parencite{GarciaLopez2020Airbnb}: limitarla podría aliviar la presión allí donde se concentra, aunque el efecto en el agregado sería menor.
\item \emph{Límites al crédito (préstamo sobre el valor o sobre la renta)}: frenan el crédito a los hogares de forma significativa; su efecto sobre los precios es más moderado \parencite{Alam2019Macroprudencial}. El FMI los recomienda para España, sobre todo para contener riesgos financieros \fmi{25-26}.
\item \emph{Compras de no residentes}: según el Banco de España, esta demanda contribuye a la subida de los precios y se concentra en Illes Balears, la Comunitat Valenciana y las Islas Canarias, donde supuso entre el 20 y el 25~\% de las compraventas en 2023 \parencite[pp.~246-247]{BdE2024Vivienda}. Una medida sobre ella actuaría, por tanto, sobre todo en esas zonas; su efecto en los precios dependería del diseño y no se cuantifica aquí.
\end{itemize}

\que{La localización.} El problema no es igual en todas partes: los problemas de acceso se concentran en las áreas urbanas y turísticas, y la creación de hogares, en las comunidades más pobladas y las grandes ciudades \parencite[pp.~232 y 242-243]{BdE2024Vivienda}. Si en unas falta vivienda y en otras hay vivienda vacía, las medidas eficaces serían locales (suelo y licencias donde falta; fiscalidad de la vacía donde abunda), y una media nacional escondería el problema. El FMI añade que la presión se ha extendido más allá de las grandes ciudades y la costa \fmi{8}.

\que{Error típico.} Recomendar «más ayudas a la compra para los jóvenes» con una oferta rígida: con los datos anteriores, buena parte acabaría en el precio. Y el contrario: descartar toda medida de demanda porque «el problema es la oferta».
\end{complemento}
